\section{Conclusion}

This paper has investigated an adaptive neural network-based control
strategy for robotic manipulators subject to modeling uncertainties,
external disturbances, and unmodeled nonlinear dynamics. By decomposing
the robot dynamics into a nominal model and a lumped uncertainty term, an
RBF neural network was introduced to approximate the unknown nonlinear
function online. Based on the nominal inverse dynamics and neural
compensation, a modular controller structure was developed, where the
trajectory tracking objective and uncertainty compensation mechanism are
separately designed.

A Lyapunov-based adaptive law was derived to update the neural network
weights, and a $\sigma$-modification technique was incorporated to improve
the robustness against persistent approximation errors and disturbances.
The stability analysis demonstrated that the closed-loop system is
uniformly ultimately bounded, and all system states and neural network
weight estimation errors remain bounded under the proposed adaptive
control framework.

Simulation studies were conducted on both a nonlinear servo system and a
two-link robotic manipulator. The results of the servo system verified
that the RBF neural network can effectively approximate unknown
non-smooth disturbances and improve tracking performance compared with the
nominal controller without compensation. Furthermore, increasing the
number of hidden nodes enhanced the approximation capability of the neural
network. The two-link manipulator simulation further demonstrated the
effectiveness of the proposed controller in the presence of dynamic
coupling, parameter uncertainties, and external disturbances, achieving
accurate trajectory tracking while maintaining bounded control inputs.

Future work will focus on experimental validation using a physical robotic
platform and further improvements of the adaptive compensation mechanism
under more challenging conditions, such as time-varying uncertainties,
actuator constraints, and communication limitations in networked robotic
systems.

